\section*{第 \ref{ch:g5:weather-and-climate} 章——天气与气候}

\begin{solution}{exo:g5:weather-and-climate:1}
(a) \omterm{def:g1:watching-the-weather:weather}{天气}（仅限今天上午）。(b) \omterm{def:g5:weather-and-climate:climate}{气候}（每年夏天）。(c) \omterm{def:g1:watching-the-weather:weather}{天气}
（昨天下午）。(d) \omterm{def:g5:weather-and-climate:climate}{气候}（每个月，年复一年）。
\end{solution}

\begin{solution}{exo:g5:weather-and-climate:2}
某地的\omterm{def:g5:weather-and-climate:climate}{气候}，是指该地多年间逐月通常出现的\omterm{def:g1:watching-the-weather:weather}{天气}。
气象学家对 30 年的测量数据取平均值。
\end{solution}

\begin{solution}{exo:g5:weather-and-climate:3}
最冷的月份是 12 月（约 6~$^\circ$C；1 月
几乎同样冷）。最湿的月份是 11 月（约 160~毫米）。否：曲线从未
低于 0~$^\circ$C。
\end{solution}

\begin{solution}{exo:g5:weather-and-climate:4}
7 月是最暖的月份，约 5~$^\circ$C。没有哪个月比
10~$^\circ$C 更暖。
\end{solution}

\begin{solution}{exo:g5:weather-and-climate:5}
\omterm{def:g3:continents-and-oceans:equator}{赤道}：热带。\omterm{def:g3:continents-and-oceans:poles}{北极}：寒带（\omterm{def:g3:continents-and-oceans:equator}{北
半球}的那一个）。\omterm{def:g3:continents-and-oceans:equator}{赤道}与\omterm{def:g3:continents-and-oceans:poles}{南极}之间的中点：\omterm{def:g3:continents-and-oceans:equator}{南半球}
温带（位于南\omterm{def:g5:weather-and-climate:tropics}{回归线}与\omterm{def:g3:continents-and-oceans:poles}{南极}
圈之间）。
\end{solution}

\begin{solution}{exo:g5:weather-and-climate:6}
希洛：热带。凤凰城和西雅图：北温带。
乌特恰维克：北寒带，位于\omterm{def:g3:continents-and-oceans:poles}{北极}圈以北。
\end{solution}

\begin{solution}{exo:g5:weather-and-climate:7}
全年：$147 + 96 + 106 + 81 + 48 + 37 + 15 + 25 + 41 + 99 + 160 + 145
= 1\,000$ 毫米。12 月、1 月和 2 月：
$145 + 147 + 96 = 388$ 毫米，这三个月的降雨量超过
全年降雨量的三分之一。
\end{solution}

\begin{solution}{exo:g5:weather-and-climate:8}
凤凰城：$35 - 14 = 21$ 度。希洛：$25 - 22 = 3$ 度。季节
对凤凰城气温的影响大得多；在希洛
则全年几乎相同。
\end{solution}

\begin{solution}{exo:g5:weather-and-climate:9}
\omterm{def:g5:weather-and-climate:rainforest}{雨林}\omterm{def:g5:weather-and-climate:climate}{气候}：全年炎热（27 或 28~$^\circ$C，没有
冬季），每个月都有大量降雨（至少 160~毫米，
全年超过 2\,000~毫米）。这座城镇最可能位于热带，靠近
\omterm{def:g3:continents-and-oceans:equator}{赤道}。
\end{solution}

\begin{solution}{exo:g5:weather-and-climate:10}
第一种解释：海滩位于热带，全年温暖，
而积雪的花园位于温带或寒带，1 月是
冬季。第二种解释：海滩位于地球\omterm{def:g3:continents-and-oceans:equator}{南半球}，
1 月是夏季，而积雪的花园位于\omterm{def:g3:continents-and-oceans:equator}{北半球}，
1 月是冬季：地球两个\omterm{def:g3:continents-and-oceans:equator}{半球}的季节
恰好相反。
\end{solution}

\begin{solution}{exo:g5:weather-and-climate:11}
气温：凤凰城要热得多，7 月约 35~$^\circ$C，
而西雅图 8 月约 20~$^\circ$C；凤凰城的 1 月（约
14~$^\circ$C）已与西雅图的 5 月一样暖和。降雨：每年不足 200~毫米，
凤凰城是\omterm{def:g5:weather-and-climate:desert}{沙漠}，西雅图约为 1\,000~毫米。因此两地的
\omterm{def:g5:weather-and-climate:climate}{气候}不同。\omterm{def:g5:weather-and-climate:climate-zone}{气候带}大致告诉我们太阳从\omterm{def:g3:continents-and-oceans:equator}{赤道}到两极
能给某地多少热量，但它们没有说明降雨多少，
也没有说明\omterm{def:g2:water-on-earth:ocean}{海洋}、\omterm{def:g4:world-physical-map:range}{山脉}或风的影响。
\end{solution}
